\begin{otherlanguage}{english}
    \begin{abstract}
    \thispagestyle{empty}
This work investigates the feasibility of using facial recognition based on \textit{embeddings} as an alternative to \textit{Near Field Communication} (\acs{NFC}) cards used for user identification in institutional environments. Initially, a conceptual and experimental validation of the approach was conducted using pre-trained models for facial feature extraction and similarity metrics in the vector space. This stage employed a synthetic facial image dataset, adopted in compliance with ethical requirements and the Brazilian General Data Protection Law (\acs{LGPD}), allowing performance evaluation through the biometric metrics \acs{FAR} (\textit{False Acceptance Rate}) and \acs{FRR} (\textit{False Rejection Rate}).

The obtained results demonstrated a clear separation between distinct identities and consistency among samples belonging to the same identity, indicating the technical feasibility of the proposed solution. Based on these findings, a functional real-time facial recognition prototype was developed using a webcam, local biometric data storage, and a graphical interface for user registration, management, and identification. The system employs the \acs{MTCNN} model for face detection and a VGGFace2 pre-trained InceptionResnetV1 model for \textit{embedding} generation. Furthermore, multi-pose enrollment, capture quality assessment, and temporal voting mechanisms were implemented to improve recognition robustness.

The results show that the proposed approach is technically viable as an alternative to physical cards, combining experimental evidence and practical validation through a functional prototype capable of performing automated and contactless biometric identification.
\\[0.5em]
\textbf{Keywords:} facial recognition; biometrics; \textit{embeddings}; computer vision; access control; \acs{NFC}.
\end{abstract}
\end{otherlanguage}
